\chapter{Adaptive Measurement Noise Matrix}
\label{ch:adaptive_r}

\TODO{}

% ---------------------------------------------------------------------
% Draft moved from Section 2.6 (Ground Truth Reconstruction): relative quantities,
% to be introduced together with the measurement error in the relative frame.
%
% Let $\mathbf{p}_e = [x_e,\,y_e]^\top$, $\psi_e$ and $v_e$ be the position, heading and
% speed of the ego vehicle, and $\mathbf{p}_o = [x_o,\,y_o]^\top$, $\psi_o$ and $v_o$
% those of the opponent, all expressed in the map frame. The position of the opponent in
% the ego frame and the relative heading are
% \begin{equation}
%     \begin{bmatrix} x_{rel} \\ y_{rel} \end{bmatrix}
%     = \mathbf{R}(\psi_e)^{\top}\,(\mathbf{p}_o - \mathbf{p}_e),
%     \qquad
%     \psi_{rel} = \psi_o - \psi_e,
% \end{equation}
% where $\mathbf{R}(\psi_e)$ is the planar rotation matrix of angle $\psi_e$. The range
% and the range rate, which are the quantities measured by the RaDAR, are computed as
% \begin{equation}
%     \rho = \|\mathbf{p}_o - \mathbf{p}_e\|,
%     \qquad
%     \dot\rho = \frac{(\mathbf{p}_o - \mathbf{p}_e)^{\top}(\mathbf{v}_o - \mathbf{v}_e)}{\rho},
% \end{equation}
% where the velocities are assumed to be aligned with the heading of each vehicle,
% $\mathbf{v}_i = v_i\,[\cos\psi_i,\,\sin\psi_i]^\top$.


\FloatBarrier

% =====================================================================
\section{Ground Truth Reconstruction}
\label{sec:gt}

The evaluation of the measurements and of the target tracking algorithm requires the
true state of the opponent at every instant. Only with this reference it is possible to
quantify the error of the raw detections of each sensor and the error of the estimates
produced by the filter, separately for each component of the state.

The reference for the opponent state is the localization of the opponent vehicle
itself. During the sessions, each team records the output of its own state estimation,
based on RTK-GNSS receivers and inertial units, and shares it after the session. The
data provided by the opponent contain, for each sample, the timestamp, the position and
the heading of the vehicle, its longitudinal speed and, when available, its
longitudinal acceleration, and they are converted into a common structure compatible
with the logs of the ego vehicle. The conversion consists in expressing the position
and the heading of the opponent in the map frame of the ego vehicle; in associating each sample of the opponent with the
sample of the ego vehicle with the closest timestamp; and in low-pass filtering the raw
inertial measurement of the longitudinal acceleration or, when it is not available, in
deriving the acceleration from the speed.

It is important to note that this reference is not exact: it is the output of the
state estimation of the opponent, and it may therefore contain errors related to the
quality of its localization. Since the errors of the sensors are computed with respect
to this reference, the analyses of the measurement errors carried out in
Chapter~\ref{ch:adaptive_r} are also affected by the accuracy of the ground truth, and part of what is observed as
an error of the detections may originate from the reference itself.
